\par\medskip\noindent\textbf{Cómo leer la salida.}
{\small\setlength{\tabcolsep}{4pt}\renewcommand{\arraystretch}{1.15}
\begin{xltabular}{\linewidth}{@{}>{\raggedright\arraybackslash}p{0.24\linewidth}>{\raggedleft\arraybackslash}p{0.15\linewidth}>{\raggedright\arraybackslash}X@{}}
\toprule Indicador & Valor & Interpretación \\ \midrule \endfirsthead
\toprule Indicador & Valor & Interpretación \\ \midrule \endhead
\bottomrule \endlastfoot
Día 0 y ventana de estimación & 2018-07-26; jul-2017 a jul-2018 & Primer día de negociación después del anuncio y período (que termina antes del evento) con que se estima el rendimiento normal. \\
Rendimiento de FB y del S\&P 500 & $-18{,}96$\,\%; $-0{,}30$\,\% & Facebook cayó casi 19\,\% mientras el mercado apenas se movió. \\
\texttt{alfa}, \texttt{beta} & 0,021; 1,279 (mercado) & Parámetros del modelo de rendimiento normal; la acción amplifica en 28\,\% los movimientos del mercado. \\
\texttt{s} & 1,34 & Desviación estándar diaria del residuo en la estimación (\%): medida de lo que es un movimiento ``normal'' propio. \\
\texttt{AR(0)} & $-18{,}59$ & Rendimiento anormal del día 0: la caída no explicada por el mercado. \\
\texttt{ee(AR0)}, \texttt{t(AR0)} & 1,34; $-13{,}85$ & El efecto es casi 14 veces su error estándar: inequívocamente significativo. \\
\texttt{CAR(-1,+1)}, \texttt{CAR(0,+5)} & $-18{,}42$; $-18{,}11$ & Rendimiento anormal acumulado: la caída no se revierte en los días siguientes; el mercado incorporó la noticia de golpe. \\
\texttt{CAR(-5,+5)} & $-15{,}76$ & Incluye días previos con rendimientos anormales levemente positivos. \\
Modelos media / ajustado & $-19{,}05$ / $-18{,}66$ & Los tres modelos dan casi el mismo resultado: con un efecto tan grande, la elección del modelo no importa. \\
Tabla por día: \texttt{AR}, \texttt{t} & solo $\tau=0$ significativo & Ningún otro día supera el valor crítico 1,97: no hubo filtración previa ni reacción tardía significativa. \\
Regresión con ficticia & $\delta=-18{,}59$; $t=-13{,}85$ & Estimar el efecto con una variable ficticia del día 0 da exactamente el mismo AR y error estándar (equivalencia del texto). \\
Excluyendo marzo de 2018 & AR(0) $-18{,}63$; $t=-14{,}47$ & Quitar de la estimación el episodio de Cambridge Analytica casi no cambia el resultado. \\
\end{xltabular}}
